\section{Additional Related work}
% \new{
We now examine related work on side-channel attacks concerning \gls{sevsnp} \glspl{cvm}, along with side-channel mitigation approaches that were not discussed in
\cref{sec:background:mitigations}.
% }
% We describe the side-channel defense landscape in this section, focusing on the paging-based \emph{side-channel}~\cite{xu2015controlledchannel} and the recent side-channel~\cite{li2021cipherleaks,li2022systematic}, since they are the most immediate threat to \glspl{cvm}.

% \{Side-channel mitigations}
% \thename is a system that aims to \emph{prevent} controlled-channel attacks from leaking sensitive information
% In this subsection, we discuss other control-channel mitigation schemes that were proposed previously, including solutions for Intel SGX.
% We choose not to mention approaches that relies



% \thename takes 

% \thename builds atop these work with an adaptive obfuscation scheme.
% Moreover, it covers full system memory accesses and does not require compiler instrumentation.

% The runtime libraries \emph{intercepts} the memory access, perform obfuscated memory accesses, and then finally \emph{translate} the memory.
% Due to the similarities between instrumentation and the real hardware MMU, these approaches are sometimes referred to as \emph{software-based MMU} schemes.
% to the runtime libraries right before an application's memory access.
% Raccon ~\cite{rane2015raccoon} \cite{borrello2021constantine}
% \cite{zhang2020klotski}
% \cite{orenbach2019cosmix}
% \cite{obelix}


% There have been proposals for obfuscated execution systems to protect SGX enclaves~\cite{ahmad2019obfuscuro,zhang2020klotski}.

% Obfuscuro and Klotski implemented obfuscated execution engines for SGX programs. Obfuscuro employs a side-channel re

% Unlike Obfuscuro and Klotski, \thename relies on the page fault handling capability.


\iffalse
	\hdr{Secure monitor for paging of \glspl{cvm}}
	Industries and academia have explored moving certain sensitive hypervisor operations (e.g., maintenance of a \gls{cvm}'s) to be performed by a privileged \emph{secure monitor} that implements security policies.
	As described in \cref{sec:discussion:cvms}, ARM CCA and Intel TDX require hypervisor paging to be delegated to a hardware-isolated secure monitor.
	For  \gls{sevsnp}, Qin et al. \cite{qin2020understanding} proposed a \gls{npf} prevention scheme that leverages a \emph{nested kernel}~\cite{dautenhahn2015nested} design to introduce a privileged monitor that executes alongside the hypervisor in host ring-0.
	This approach requires intrusive modifications to the host hypervisor.
	% which is impractical for cloud computing scenarios that efficient lifting and shifting of workloads~\cite{ge2022hecate}.

	A challenge with monitor-based defenses is to define a security policy that does not hamper the hypervisor's existing capabilities.
	Strict policies like preventing \gls{cvm}'s memory pages from being unmapped would severely limit the flexibility in memory management and thus are undesirable.
	As a consequence, we expect that primitives that enable controlled-channel attacks will eventually emerge as with the case of \emph{page blocking} mechanism of Intel \gls{tdx}~\cref{sec:discussion:cvms}.
\fi

\subsection{Side-channel attacks on SEV-SNP}
\label{subsec:related-work:attacks}
% Previous attacks on SEV-protected CVMs that utilize side channels (\cref{sec:background:side-channels}) as a key component of their strategy.
The use of side-channel in previous attacks can be categorized into two main types: (1) identifying sensitive pages or function executions as a preliminary step before initiating more targeted attacks~\cite{li2021cipherleaks,li2022systematic,li2019exploiting,cachewarp,sevstep,werner2019severest,morbitzer2018severed,heckler,wesee}, and (2) extracting secrets from the CVM's memory accesses that depend on sensitive information ~\cite{li2021cipherleaks,li2022systematic,sevstep,ciphersteal}.
While most of the SEV-specific vulnerabilities are already addressed in the latest version of SEV-SNP, e.g., \cite{li2021cipherleaks,cachewarp,li2019exploiting,werner2019severest,morbitzer2018severed}, their side-channel components remain difficult to mitigate.
\iffalse
	\thename's adaptive rerandomization scheme targets the \emph{temporal reslution} of the attacks, rendering future attacks much more challenging.
\fi

\subsubsection{Profiling using controlled channel}
\label{sec:relatedwork:profiling}
All existing attacks on SEV and its extensions require a \emph{profiling} phase with the coarse-grained \gls{npf} controlled channel (\cref{sec:background:controlled-channel}) to infer locations of sensitive code/data pages within the guest's \gls{gpa} space.
Profiling attacks often undermine KASLR through statistical analysis~\cite{heckler,wesee,morbitzer2018severed,li2022systematic} or by matching page access patterns with known ones~\cite{li2019exploiting,wilke2020sevurity}.
\iffalse
	We found that most of these attacks necessitate the accumulation of page access traces that reach a specific critical mass, e.g., $22,400$ \gls{cvm} memory accesses since boot~\cite{heckler}, or all page accesses and their instruction counts during an HTTPS request~\cite{li2022systematic}.
	As shown in our evaluations (\cref{subsec:profiling-attack}), profiling attempts will inevitably be met with constant rerandomizations.
\fi

% Essentially, profiling attacks attempt to break KALSR with statistical analysis~\cite{heckler,wesee,morbitzer2018severed,li2022systematic} and page access pattern matching~\cite{li2019exploiting,wilke2020sevurity}.
% Most attacks require an accumulation of page traces that reach a certain critical mass, e.g., 22,400 memory accesses from boot to identify the OpenSSH pages~\cite{heckler}, or page accesses and their instruction counts throughout a HTTPS request~\cite{li2022systematic}.
% If a \thename address space rerandomization is encountered during profiling, the attack is disrupted before yielding fruitful results.


% This information could be obtained through assumed posterior knowledge of the target CVM execution~\cite{li2021cipherleaks,cachewarp,werner2019severest},


\hdr{CVM secret extraction}
\label{sec:relatedwork:extraction}
Secret extraction attacks require prior identification of sensitive code and data, allowing the extraction of secret-dependent memory accesses.
Previous work discussed secret extraction through fine-grained side channels~\cite{ciphersteal,li2022systematic,li2021cipherleaks,sevstep} whose temporal resolution is maximized through the single-stepping technique (\cref{sec:background:interrupts}).
SEV-Step~\cite{sevstep} presents an L2 cache Prime+Probe attack on the AES T-table for key extraction.
Li and Wilke et al.~\cite{li2022systematic} enhance the now-patched \texttt{CipherLeak}~\cite{li2021cipherleaks} ciphertext attack to perform secret key extraction without relying on the encrypted VMSA.
\texttt{CipherSteal}~\cite{ciphersteal} infers input data in protected DNN computation using the ciphertext side channel.
% Aside from introducing difficulties in locating sensitive code and data pages,
% Although not designed to mitigate attacks using fine-grained side channels completely, 
\iffalse
	\thename would increase the rate of rerandomization during these attacks.
	Each rerandomization effectively pushes the attacker back to the profiling stage.
\fi
% disrupt the attack before they yield fruitful results.

% As another example, the \texttt{Prime+Probe} attack against L2 cache on AMD SEV~\cite{sevstep} requires tracing access to the AES T-table from 70 XTS decryption operations using single-stepping (\cref{sec:background:side-channels}), which requires at least $5,600$ cache access traces (70 $\times$ 80 accesses to the T-table per decryption). 
% The two fine-grained side channels that are still prevalent on SEV are the ciphertext side-channel and L2 cache



% Attacks on cryptographic implementations require tracking the page locations of cryptographic functions and sensitive buffers\cite{li2021cipherleaks,li2022systematic,sevstep}.
% CacheWrap~\cite{cachewarp}, CipherLeak~\cite{li2021cipherleaks} requires identification of the starting point of the target function of the attack.



% show how the side-channel attack primitives are chained to form a practical attack: a coarse-grained primitive is used for the \emph{profiling} phase while fine-grained ones are used for the \emph{extraction} phase.

% In the \emph{profiling phase}, the adversary aims to pinpoint the sensitive pages/function executions within the \gls{gpa} space of the \gls{cvm}.
% The adversary resorts to the coarse-grained \gls{npf} controlled channel that provides more bandwidth in this phase.

\iffalse
	Since the profiling phase has to be performed on the entire memory space of the \gls{cvm},
	the fast \gls{npf} controlled-channel attack is used as a practical primitive to deterministically trace the page-level memory accesses of the \gls{cvm}~\cite{li2021cipherleaks,li2022systematic,li2019exploiting,cachewarp,sevstep,werner2019severest,morbitzer2018severed,heckler,wesee}.


	Then, the adversary proceeds to the \emph{extraction phase} in which finer-grained attack primitives, including the finer-grained side channel primitives, are used.
	The extraction phase focuses on recovering secrets from the \gls{cvm} by utilizing the information obtained from secret-dependent memory accesses exposed through side channels.
	This phase often depends on the exact locations of the target code or data collected in the profiling phase, for example, the exact pages that store AES T-table~\cite{sevstep}.
	With the targeted page locations established, the attacker can utilize lower-bandwidth yet high-precision side channels, such as cache and ciphertext, to carry out the extraction~\cite{sevstep,li2021cipherleaks}.

	\hjcomm{Don't we need this?}{...}
	~\cite{li2022systematic,werner2019severest,morbitzer2018severed,li2021cipherleaks,heckler,wesee,sevstep}.
\fi


\subsection{Approaches for side-channel mitigation}
Aside from the approaches we discussed in \cref{sec:background:mitigations}, there are other proposals for side-channel mitigation, but with certain limitations.
Raccoon~\cite{rane2015raccoon} is the first to discuss obfuscation for side-channel mitigation in the context of SGX by incorporating decoy execution and ORAM, which incurs up to 1000$\times$ overheads on simple programs.
Shinde et al.~\cite{shinde2016preventing} propose a deterministic multiplexing that forces sensitive code and data accesses onto a single page but incurs up to 4000$\times$ overheads on program execution.
SGX-Shield~\cite{seo2017sgxshield} introduces load-time address space randomization for enclaves but would be defeated through online access profiling.
OBLIVIATE~\cite{ahmad2018obliviate} introduces an obfuscated file system for SGX enclaves.
DR.SGX~\cite{brasser2019drsgx} proposes periodic runtime enclave memory rerandomization at the cache-line level.
ZeroTrace~\cite{sasy2018zerotrace} implements a memory interface that uses ORAM to obfuscate memory accesses.
Aside from introducing significant overheads, DR.SGX and ZeroTrace only support the obfuscation of data accesses.
Finally, it is unclear how previous systems would be adapted to support CVM full-system memory access obfuscation.
% Qin et al.~\cite{qin2023protecting} proposed refactoring the hypervisor to isolate access to paging through a \emph{Virtualization Security Module (VSM)} to mitigate the \gls{npf} controlled channel.
% This approach necessitates maintaining complex invariants to safeguard the VSM from the untrusted hypervisor~\cite{dautenhahn2015nested} and requires intrusive modifications to the hypervisor codebase, which is not ideal in cloud computing.

% Additionally, it requires establishing parts of the hypervisor as the TCB through the secure boot.
% which may complicates VM migrations across different hosts.
% For the control channel mitigation of \gls{sevsnp}


% It also must tackle unique challenges in retrofitting, for instance, rerandomizing whole-system pages, including the page table itself.


% Several \gls{tee} prototypes such as Keystone~\cite{keystone} and Sanctum~\cite{costan2016sanctum} employ a dual page table approach: the sensitive enclave pages are mapped in a dedicated page table, the \emph{enclave page table}, that are hardware-isolated from the untrusted OS.
% In these systems, the enclaves self-handle their page faults, removing the controlled channel from the OS. 
% This approach is limited as it is not compatible with the current OS memory management model, which uses demand-paging.
% Extending these \gls{tee} designs, InvisiPage~\cite{aga2019invisipage} and Autarky~\cite{orenbach2020autarky} proposed ISA extensions to Intel SGX to enable demand-paging by facilitating cooperation between the enclave and the untrusted OS for ORAM-based page swapping.

% Unlike the mentioned systems that require customized CPU designs or hardware modifications, \thename can be implemented on commodity \gls{sevsnp} hardware and other \gls{cvm} architectures.
% Retrofitting an adaptive obfuscation scheme into the trust boundary of \gls{cvm} also poses unique challenges that \thename must address, namely the continuous randomization of whole-system pages, including the page table. 

% Additionally, without hardware-enforced page table isolation, \thename's paging scheme achieves \gls{pao} for both data and page table pages by integrating a recursive ORAM scheme into the multi-level page table, thus ensuring comprehensiveness.

% integrate a recursive ORAM scheme into the multi-level page table.
% While InvisPage assumes hardware support for enclave-OS page table separation, Autarky only introduces minimal hardware changes that (1) hide the faulting address from the OS and (2) force the OS to notify the enclave about page faults.
% Both works proposed using ORAM to swap EPC pages in and out of the enclave without leaking the access pattern.


% Since \thename relies on the natural isolation of the guest page table in \glspl{cvm}, it does not need new hardware extensions like the previous work.

% Moreover, \thename does not rely on compiler instrumentation like Autarky.



% InvisiPage requires intrusive hardware modifications to support enclave-private page tables. On the other hand, Autarky ...



% \subsection{Dynamic program instrumentation}
% \label{sec:rw:dyn-inst}
